\chapter{Hiérarchie de la mémoire et caches}\label{ch:ll:memory-hierarchy-and-caches}

Sur le portable de ce livre, la lecture d'un niveau de prix d'un carnet d'ordres coûte environ une nanoseconde lorsque le niveau se trouve dans le
cache de premier niveau du cœur, trois ou quatre nanosecondes lorsqu'il se trouve au deuxième niveau, et environ 150~nanosecondes lorsqu'il
doit être lu en mémoire~: un facteur supérieur à cent pour une même instruction. Le carnet jouet du chapitre~1, bâti sur des arbres équilibrés,
dont les nœuds sont dispersés dans le tas, passait l'essentiel de son temps à attendre ce genre de lecture. La plus grande part de l'écart
entre un carnet d'ordres rapide et un carnet lent ne tient pas à l'algorithme mais à l'endroit où vivent ses octets, et ce chapitre porte sur cela~:
les niveaux de la hiérarchie et leur coût, ce qui se passe lorsque deux cœurs partagent une ligne, la manière dont les adresses virtuelles
sont traduites, et ce que le matériel fait pour masquer tout cela.

\section{Niveaux de cache, lignes et associativité}

\begin{definition}[Ligne de cache, cache associatif par ensembles]\label{def:ll:memory-hierarchy-and-caches:line}
Les données circulent entre les niveaux de la hiérarchie par blocs de taille fixe appelés \emph{lignes de cache}\index{ligne de cache} (cache lines), de 64~octets
sur les processeurs x86 actuels. Dans un \emph{cache associatif par ensembles}\index{cache associatif par ensembles} (set-associative cache) à $S$ ensembles de $W$ voies, une
ligne ne peut être stockée que dans l'ensemble désigné par les bits de son adresse, dans l'un quelconque des $W$ emplacements de cet ensemble~; le cache contient
$S \cdot W$ lignes.
\end{definition}

\begin{definition}[Défaut de cache, ensemble de travail]\label{def:ll:memory-hierarchy-and-caches:miss}
Un \emph{défaut de cache}\index{défaut de cache} (cache miss) est un accès à une ligne absente d'un niveau donné, qu'il faut alors
aller chercher au niveau inférieur. L'\emph{ensemble de travail}\index{ensemble de travail} (working set) d'un morceau de code sur un intervalle est l'ensemble des
lignes qu'il touche sur cet intervalle~; le niveau dont la capacité peut l'accueillir fixe sa vitesse.
\end{definition}

Une adresse se décompose en trois champs (\cref{fig:ll:memory-hierarchy-and-caches:split})~: les 6 bits de poids faible donnent le décalage
dans la ligne de 64~octets~; les $\log_2 S$ bits suivants choisissent l'ensemble~; le reste, l'étiquette, identifie la ligne parmi celles qui
correspondent à l'ensemble. Deux conséquences pour l'organisation des données. Une structure à cheval sur deux lignes coûte deux défauts là où un seul
suffirait. Et des adresses qui diffèrent d'un multiple de $S \times 64$ octets se disputent les mêmes $W$ emplacements~: un programme qui parcourt
plusieurs tableaux dont les pas sont ce multiple (une puissance de deux, typiquement 4\,KiB) peut subir des défauts dans un cache
presque vide.

\begin{omfigure}
\begin{tikzpicture}[font=\footnotesize]
\draw[fill=gray!10] (0,0) rectangle (6.4,0.6); \node at (3.2,0.3) {étiquette};
\draw[fill=omDef!15] (6.4,0) rectangle (8.8,0.6); \node[font=\scriptsize] at (7.6,0.3) {indice d'ensemble};
\draw[fill=omThm!12] (8.8,0) rectangle (10.8,0.6); \node at (9.8,0.3) {décalage};
\node[anchor=south,align=center] at (3.2,0.65) {bits restants de l'adresse};
\node[anchor=south,align=center] at (7.6,0.65) {bits 6--11\\ (64 ensembles)};
\node[anchor=south,align=center] at (9.8,0.65) {bits 0--5\\ (64 octets)};
\draw[-Stealth,thick,omDef] (7.6,0) -- (7.6,-0.75);
\foreach \w in {0,...,11} \draw[draw=omDef,fill=omDef!6] (4.6+0.5*\w,-0.8) rectangle (5.05+0.5*\w,-1.3);
\node[anchor=north,align=center] at (7.35,-1.4) {l'ensemble choisi~: 12 voies, dont chacune peut contenir une ligne portant cet indice};
\end{tikzpicture}

\omcaption{Comment une adresse désigne son emplacement dans un \omterm{def:ll:memory-hierarchy-and-caches:line}{cache associatif par ensembles}. Un cache de 48~KiB à 12 voies, aux lignes de 64~octets, compte 64
ensembles~; les bits 6 à 11 de l'adresse désignent l'ensemble, si bien que deux adresses distantes de 4\,KiB partagent un même ensemble.}\label{fig:ll:memory-hierarchy-and-caches:split}
\end{omfigure}

\begin{example}[Géométrie d'un cache de premier niveau]\label{ex:ll:memory-hierarchy-and-caches:geometry}
Un cache de 48~KiB à 12 voies, aux lignes de 64~octets, compte $48 \times 1024 / (64 \times 12) = 64$ ensembles. L'indice d'ensemble est formé des bits 6 à
11~; des adresses égales modulo $64 \times 64 = 4\,096$ octets tombent dans le même ensemble. Treize structures très sollicitées placées
au même décalage de treize pages de 4~KiB ne peuvent toutes rester dans ce cache, alors qu'elles n'occupent ensemble que 832 octets.
\end{example}

La \Cref{fig:ll:memory-hierarchy-and-caches:stairs} mesure la hiérarchie par un parcours de pointeurs~: un tampon de $n$ emplacements, un par
\omterm{def:ll:memory-hierarchy-and-caches:line}{ligne de cache}, disposé en un unique cycle aléatoire, chaque emplacement contenant l'adresse du suivant. L'adresse de chaque lecture est
le résultat de la lecture précédente, si bien que deux lectures ne se recouvrent jamais et que le temps par étape est la latence du niveau dans lequel le
tampon tient. L'escalier montre le premier niveau jusqu'à 48~KiB, à environ \qty{1.1}{\nano\second} (cinq cycles), le deuxième
niveau jusqu'à environ un mégaoctet et demi, à environ \qty{3.5}{\nano\second}, le troisième niveau, d'environ trois à huit
mégaoctets, entre 15 et \qty{23}{\nano\second}, proche de sa latence publiée de 75 cycles ou plus (environ
\qty{20}{\nano\second}), et la mémoire à partir d'environ douze mégaoctets, entre 100 et \qty{175}{\nano\second}~: la
mémoire basse consommation d'un portable, qu'une mesure publiée situe entre 150 et 175 nanosecondes pour ce processeur. Le
troisième niveau fait 24~MiB, partagé par tous les cœurs~; un parcours aléatoire dans les pages d'une machine virtuelle commence à manquer dans ce
niveau bien avant sa taille nominale.

\begin{omfigure}
\begin{tikzpicture}
\begin{axis}[width=12cm,height=6cm,xmode=log,ymode=log,log basis x=2,xmin=4096,xmax=268435456,ymin=0.8,ymax=500,
  xtick={4096,32768,262144,2097152,16777216,134217728},xticklabels={4 KiB,32 KiB,256 KiB,2 MiB,16 MiB,128 MiB},
  ytick={1,2,5,10,20,50,100,200},yticklabels={1,2,5,10,20,50,100,200},
  xlabel={ensemble de travail (échelle log)},ylabel={ns par lecture (échelle log)},
  tick label style={font=\small},label style={font=\small},grid=major,grid style={gray!25},
  legend columns=3,legend style={font=\footnotesize,at={(0.5,-0.3)},anchor=north,draw=none,fill=none,
  /tikz/every even column/.append style={column sep=8pt}},table/col sep=comma]
\addplot[omDef,very thick,mark=*,mark size=1.2pt] table[x=bytes,y=random_ns]{figdata/low-latency/03-memory-hierarchy-and-caches/measured_stairs.csv};
\addplot[omThm,thick,mark=square*,mark size=1.2pt] table[x=bytes,y=random_huge_ns]{figdata/low-latency/03-memory-hierarchy-and-caches/measured_stairs.csv};
\addplot[omMeth,thick,dashed,mark=triangle*,mark size=1.2pt] table[x=bytes,y=sequential_ns]{figdata/low-latency/03-memory-hierarchy-and-caches/measured_stairs.csv};
\legend{ordre aléatoire,{ordre aléatoire, grandes pages},ordre des adresses}
\end{axis}
\end{tikzpicture}

\omcaption{L'escalier de la mémoire~: temps par lecture dépendante à travers un tampon d'un emplacement par \omterm{def:ll:memory-hierarchy-and-caches:line}{ligne de cache}, en fonction de la
taille du tampon. Dans l'ordre des adresses, le \omterm{def:ll:memory-hierarchy-and-caches:prefetch}{préchargement matériel} masque presque tout. Mesuré sur un portable (Intel Core Ultra 7 155H)
sous WSL2, sans cœurs isolés, fixé à un seul CPU, la machine étant par ailleurs inactive. Données~: \texttt{bench\_mem.py}.}\label{fig:ll:memory-hierarchy-and-caches:stairs}
\end{omfigure}

\section{Cohérence et faux partage}

\begin{definition}[Cohérence de cache]\label{def:ll:memory-hierarchy-and-caches:coherence}
La \emph{cohérence de cache}\index{cohérence de cache} (cache coherence) est le protocole par lequel les caches privés de plusieurs cœurs gardent une
vision commune de chaque ligne~: avant qu'un cœur n'écrive dans une ligne, il doit l'obtenir en état exclusif, en invalidant la copie de tous les autres
cœurs~; une lecture ultérieure par un autre cœur doit aller chercher la ligne modifiée auprès de son auteur. Dans la famille MESI, chacune des lignes
de chaque cache est modifiée, exclusive, partagée ou invalide.
\end{definition}

\begin{definition}[Faux partage]\label{def:ll:memory-hierarchy-and-caches:false}
Le \emph{faux partage}\index{faux partage} (false sharing) survient lorsque deux cœurs écrivent à plusieurs reprises des variables différentes situées dans la
même \omterm{def:ll:memory-hierarchy-and-caches:line}{ligne de cache}~: aucun des deux ne partage de données avec l'autre, et pourtant chaque écriture invalide la copie de l'autre cœur, et la ligne
circule entre eux.
\end{definition}

Un thread de données de marché qui met à jour un compteur de séquence, et un thread de stratégie qui met à jour ses propres statistiques, placés
côte à côte dans une même structure, paient le prix du déplacement de la ligne d'un cœur à l'autre à chaque mise à jour, sans que
la moindre donnée soit partagée. La \Cref{fig:ll:memory-hierarchy-and-caches:false} le mesure~: un thread qui incrémente un compteur paie environ
\qty{4}{\nano\second} par incrémentation atomique~; deux threads qui incrémentent chacun leur propre compteur sur des lignes distinctes paient environ
la même chose~; deux threads dont les compteurs partagent une même ligne paient environ cinq fois plus. Le remède tient à la disposition~: donner à chaque
variable écrite indépendamment sa propre ligne (\texttt{CachePadded}, ci-dessous), et garder ensemble les variables écrites par un même thread
à l'écart de celles qu'écrit un autre.

\begin{omfigure}
\begin{tikzpicture}
\begin{axis}[width=9cm,height=5cm,xbar,bar width=12pt,xmin=0,enlarge x limits={upper,value=0.15},
  symbolic y coords={same line,padded,one thread},ytick=data,yticklabels={un seul thread,avec remplissage,même ligne},enlarge y limits=0.3,
  xlabel={ns par incrémentation atomique},tick label style={font=\small},label style={font=\small},
  grid=major,grid style={gray!25},nodes near coords,nodes near coords style={font=\footnotesize},
  every node near coord/.append style={/pgf/number format/fixed,/pgf/number format/precision=1},
  table/col sep=comma]
\addplot[fill=omThm!45,draw=omThm] table[y=layout,x=ns_per_increment]{figdata/low-latency/03-memory-hierarchy-and-caches/measured_false.csv};
\end{axis}
\end{tikzpicture}

\omcaption{\omterm{def:ll:memory-hierarchy-and-caches:false}{Faux partage}~: deux threads épinglés sur deux CPU, incrémentant chacun son propre compteur atomique 20 millions de fois,
avec les compteurs sur une même \omterm{def:ll:memory-hierarchy-and-caches:line}{ligne de cache} ou sur deux. Mesures effectuées sur un portable (Intel Core Ultra 7 155H) sous WSL2, sans
cœurs isolés. Données~: \texttt{bench\_mem.py}.}\label{fig:ll:memory-hierarchy-and-caches:false}
\end{omfigure}

\section{Tampons de traduction d'adresses et grandes pages}

\begin{definition}[Tampon de traduction d'adresses, grande page]\label{def:ll:memory-hierarchy-and-caches:tlb}
Les programmes utilisent des adresses virtuelles, traduites en adresses physiques par les tables de pages. Le \emph{tampon de traduction
d'adresses}\index{tampon de traduction d'adresses} (translation lookaside buffer, TLB) conserve les traductions récentes~; un défaut de TLB oblige le processeur à parcourir les
tables de pages, plusieurs accès mémoire dépendants. Une \emph{grande page}\index{grande page} (huge page) est une page plus grande que la page
de base de 4~KiB (2~MiB ou 1~GiB sur x86-64), couverte par une seule entrée du TLB.
\end{definition}

La portée du TLB est son nombre d'entrées multiplié par la taille de page. Un TLB de second niveau de 2\,048 entrées couvre 8~MiB de
pages de 4~KiB et 4~GiB de pages de 2~MiB~; un carnet d'ordres, une table de symboles et quelques anneaux d'un processus de trading peuvent facilement
dépasser le premier et jamais le second. Linux propose les \omterm{def:ll:memory-hierarchy-and-caches:tlb}{grandes pages} de deux façons~: des pages réservées dans un pool, et les \omterm{def:ll:memory-hierarchy-and-caches:tlb}{grandes pages} transparentes,
que le noyau constitue de lui-même~; en mode par défaut \texttt{madvise}, il ne le fait que pour les régions que le
programme a marqué avec \texttt{madvise(MADV\_HUGEPAGE)}. La documentation du noyau note que le gain est plus important
sous virtualisation, où chaque parcours passe par deux niveaux de tables de pages, celui de l'invité et celui de l'hyperviseur. Sur
ce portable, qui est virtualisé, un parcours aléatoire sur 256~MiB coûte environ \qty{180}{\nano\second} par pas avec des pages
de 4~KiB et environ \qty{150}{\nano\second} avec les \omterm{def:ll:memory-hierarchy-and-caches:tlb}{grandes pages} (\texttt{measured\_huge.csv})~: la différence tient
au parcours des tables de pages.


\section{Préchargement et schémas d'accès}

\begin{definition}[Préchargement matériel]\label{def:ll:memory-hierarchy-and-caches:prefetch}
Le \emph{préchargement matériel}\index{préchargement matériel} (hardware prefetching) est le chargement, par le processeur, de lignes avant leur demande,
en reconnaissant des régularités dans les \omterm{def:ll:memory-hierarchy-and-caches:miss}{défauts de cache} récents~: lignes consécutives, pas constants.
\end{definition}

La courbe en tirets de la \cref{fig:ll:memory-hierarchy-and-caches:stairs} reprend le même parcours, avec les emplacements visités dans
l'ordre des adresses~: le préchargeur reconnaît le flux et les chargements coûtent moins de \qty{2}{\nano\second} tant que le tampon quitte le
second niveau, et de quatre ou cinq nanosecondes depuis la mémoire, soit une trentaine de fois moins que le parcours aléatoire. Rien n'a changé que
l'ordre. Les règles pratiques en découlent~: ranger ensemble, dans des tableaux contigus, ce que lit le chemin critique (hot path), et non dans des nœuds
liés par des pointeurs~; itérer dans l'ordre des adresses~; garder les données chaudes d'un message (prix, quantité, côté) dans la première ligne de sa
structure~; et lorsque l'accès est réellement aléatoire, comme une recherche par identifiant d'ordre, rendre la table assez petite pour rester dans le cache
(chapitre~19).

\section{Tutoriel~: mesurer l'escalier}\label{tut:ll:memory-hierarchy-and-caches:tut}

\begin{tutorial}
\textbf{Objectif.} Mesurer la latence de chaque niveau de la hiérarchie, le coût du \omterm{def:ll:memory-hierarchy-and-caches:false}{faux partage} (false sharing) et l'effet des \omterm{def:ll:memory-hierarchy-and-caches:tlb}{grandes
pages}. \textbf{État final~:} les \cref{fig:ll:memory-hierarchy-and-caches:stairs,fig:ll:memory-hierarchy-and-caches:false} et
le tableau des \omterm{def:ll:memory-hierarchy-and-caches:tlb}{grandes pages}.
\begin{tutsteps}
\item \textbf{Un anneau de pointeurs.} Mélanger l'ordre des emplacements avec Fisher--Yates, puis stocker dans chaque emplacement l'adresse de la
suivante~: un cycle unique parcourant toutes les lignes, que nul préchargeur ne peut prévoir.
\omcode{code/firm/memkit/cpp/firm_memkit.hpp}{52}{72}{Construction et parcours d'un anneau de pointeurs.}\label{lst:ll:memory-hierarchy-and-caches:chase}
\item \textbf{\omterm{def:ll:memory-hierarchy-and-caches:tlb}{Grandes pages} sur demande.} \texttt{Buffer} associe une mémoire anonyme, demande des \omterm{def:ll:memory-hierarchy-and-caches:tlb}{grandes pages} transparentes avec
\texttt{madvise}, et écrit chaque page une fois pour qu'aucun défaut de page ne survienne plus tard sur le chemin critique (chapitre~6).
\omcode{code/firm/memkit/cpp/firm_memkit.hpp}{27}{36}{Un tampon avec grandes pages et pré-chargement des pages (pre-faulting).}\label{lst:ll:memory-hierarchy-and-caches:buffer}
\item \textbf{Un compteur par ligne.} \texttt{CachePadded} aligne sa valeur sur une ligne qui lui est propre~; deux threads incrémentent
leurs compteurs par additions atomiques relâchées, une fois avec remplissage, une fois regroupés dans une même structure.
\omcode{code/firm/memkit/cpp/firm_memkit.hpp}{22}{26}{Une valeur seule sur sa ligne de cache.}\label{lst:ll:memory-hierarchy-and-caches:padded}
\item \textbf{Lancer} \texttt{python bench\_mem.py} (une minute environ)~: l'escalier avec des pages de 4~KiB et des \omterm{def:ll:memory-hierarchy-and-caches:tlb}{grandes pages}, et dans l'ordre des
adresses, la comparaison des \omterm{def:ll:memory-hierarchy-and-caches:tlb}{grandes pages} sur 256~MiB, et le test de \omterm{def:ll:memory-hierarchy-and-caches:false}{faux partage} sur les CPU 2 et 4.
\end{tutsteps}
\textbf{Pour aller plus loin.} Utiliser un pas de 4\,096 octets au lieu de 64 dans le parcours et trouver l'\omterm{def:ll:memory-hierarchy-and-caches:miss}{ensemble de travail} (working set) pour lequel
le premier niveau ne suffit plus, bien plus tôt que 48~KiB~; épingler les deux threads de \omterm{def:ll:memory-hierarchy-and-caches:false}{faux partage} sur un même CPU et expliquer le résultat.
\end{tutorial}

\section{Construction~: outils de placement en mémoire}\label{bld:ll:memory-hierarchy-and-caches:memkit}

\begin{build}
\bfield{Rôle} Les primitives de disposition qu'utilise chaque composant ultérieur~: compteurs et numéros de séquence écrits par des threads différents
(anneaux, chapitre~12~; statistiques du feed handler, chapitre~18) et tampons qui ne doivent provoquer ni défaut de page ni défaut de TLB
sur le chemin critique (carnets d'ordres, journaux).
\bfield{Interface} C++20 \texttt{firm::memkit}~: \texttt{kLine}, \texttt{CachePadded<T>},
\texttt{Buffer(bytes, huge)} avec \texttt{data()}, \texttt{size()}, \texttt{huge\_requested()},
\texttt{chase\_ring(base, n, stride, random, seed)}, \texttt{chase(start, steps)}. Rust \texttt{firm\_memkit}~: \texttt{CachePadded<T>}, \texttt{chase\_ring},
\texttt{chase} sur un tampon d'octets.
\bfield{Règles} Les valeurs avec remplissage occupent des lignes entières~; les tampons sont alignés sur les pages, arrondis à la taille de page demandée, et
pré-chargées à la construction~; une demande de \omterm{def:ll:memory-hierarchy-and-caches:tlb}{grande page} échouée se replie sur les pages de base et le signale.
\bfield{Tests d'acceptation} \texttt{code/firm/memkit/}~: tailles et alignements des valeurs avec remplissage, valeurs voisines avec
remplissage à une ligne d'écart, un tampon de \omterm{def:ll:memory-hierarchy-and-caches:tlb}{grandes pages} arrondi à 2~MiB, et l'anneau formant un cycle unique dans les deux ordres, en C++ et en Rust.
\bfield{Pour aller plus loin} Signaler les \omterm{def:ll:memory-hierarchy-and-caches:tlb}{grandes pages} effectivement obtenues, à partir de \texttt{/proc/self/smaps}~; allouer depuis des pages réservées de 1~GiB
lorsque l'hôte en dispose.
\end{build}

\begin{omsources}
\item U.~Drepper, ``What every programmer should know about memory'', LWN.net, 2007.
\item M.~Papamarcos et J.~Patel, ``A low-overhead coherence solution for multiprocessors with private cache memories'',
ISCA 1984.
\item Documentation du noyau Linux, ``Transparent hugepage support''.
\item Chips and Cheese, ``Intel's Redwood Cove: baby steps are still steps'' (2024) and ``Update on Meteor Lake DRAM
latency measurements'' (2024).
\end{omsources}

\section{Exercices}

\begin{exercise}[$\star$]\label{exo:ll:memory-hierarchy-and-caches:1}
Dans un cache dont les \omterm{def:ll:memory-hierarchy-and-caches:line}{lignes de cache} font 64~octets et qui compte 64 ensembles, à quel ensemble correspond l'adresse \texttt{0x7f3a12345e48}, et quel est son décalage
au sein de sa \omterm{def:ll:memory-hierarchy-and-caches:line}{ligne de cache}~?
\end{exercise}

\begin{exercise}[$\star$]\label{exo:ll:memory-hierarchy-and-caches:2}
Un cache de 2~MiB à 16 voies, dont les \omterm{def:ll:memory-hierarchy-and-caches:line}{lignes de cache} font 64~octets. Combien d'ensembles possède-t-il, et à quel pas les adresses tombent-elles dans le
même ensemble~?
\end{exercise}

\begin{exercise}[$\star$]\label{exo:ll:memory-hierarchy-and-caches:3}
Un processus utilise 600~MiB de données chaudes. Combien d'entrées de TLB lui faut-il avec des pages de 4~KiB, et avec des pages de 2~MiB~?
\end{exercise}

\begin{exercise}[$\star\star$]\label{exo:ll:memory-hierarchy-and-caches:4}
Une structure de message contient un numéro de séquence de 8~octets, écrit par le thread de flux, et un compteur de 4~octets
écrit par le thread de stratégie, côte à côte avec lui. Que se passe-t-il, et comment y remédier, au prix de combien d'octets supplémentaires~?
\end{exercise}

\begin{exercise}[$\star\star$]\label{exo:ll:memory-hierarchy-and-caches:5}
D'après la \cref{fig:ll:memory-hierarchy-and-caches:stairs}, combien de temps prend une mise à jour du carnet d'ordres qui effectue quatre lectures aléatoires dépendantes
lorsque ses données tiennent dans le deuxième niveau, puis lorsqu'elles proviennent de la mémoire~?
\end{exercise}

\begin{exercise}[$\star\star$]\label{exo:ll:memory-hierarchy-and-caches:6}
Pourquoi le parcours dans l'ordre des adresses reste-t-il rapide bien au-delà du deuxième niveau, alors que le parcours aléatoire ne l'est pas~? Que
cela implique-t-il pour les listes chaînées et les arbres équilibrés sur le chemin critique~?
\end{exercise}

\begin{exercise}[$\star\star\star$]\label{exo:ll:memory-hierarchy-and-caches:7}
\emph{Programmation.} Réécrire le \texttt{chase} en Rust de \texttt{firm\_memkit} avec des pointeurs bruts et \texttt{unsafe}, comme le fait le C++,
puis comparer les deux versions à la version à indices sur un anneau de 64~MiB. La vérification des bornes est-elle visible dans le temps par pas~? Pourquoi,
ou pourquoi pas~?
\end{exercise}

\begin{exercise}[$\star\star\star$]\label{exo:ll:memory-hierarchy-and-caches:8}
\emph{Trouver la faille.} «~Nous avons mesuré la mise à jour de notre carnet d'ordres à \qty{20}{\nano\second} en appliquant les mêmes messages d'ajout et
de suppression à un seul instrument dix millions de fois, dans une boucle.~»
\end{exercise}

\section{Problème~: l'escalier}

\begin{problem}[{Problème du week-end --- combien de carnets tiennent}]\label{pb:ll:memory-hierarchy-and-caches:1}
Un processus de market making cote un ensemble d'instruments. Pour chacun, il conserve une échelle de prix dense de 256 niveaux par côté,
chaque niveau occupant 32~octets (prix, quantité, nombre d'ordres, indicateurs), plus 4~KiB d'état par instrument. L'équipe veut que
les données chaudes de chaque instrument tiennent dans le cache de deuxième niveau du cœur qui exécute la stratégie.

\medskip\noindent\textbf{Partie I --- La mesure.}
\begin{enumerate}
\item Lire la latence du premier niveau, celle du deuxième niveau et celle de la mémoire à partir de
la \cref{fig:ll:memory-hierarchy-and-caches:stairs}.
\item Où se trouvent les marches, et quelles capacités suggèrent-elles~?
\item Pourquoi la marche du troisième niveau se termine-t-elle vers huit mégaoctets, et non à 24~MiB~?
\item Qu'est-ce qui fixe la hauteur de la marche entre le deuxième niveau et la mémoire~?
\end{enumerate}

\medskip\noindent\textbf{Partie II --- La disposition.}
\begin{enumerate}[resume]
\item Combien d'octets occupent les données chaudes d'un instrument~?
\item Combien de \omterm{def:ll:memory-hierarchy-and-caches:line}{lignes de cache} cela représente-t-il~?
\item Si les trois quarts d'un deuxième niveau de 2~MiB sont disponibles pour les carnets, combien d'instruments peuvent tenir~?
\item Combien en tiennent, au maximum, dans un premier niveau de 48~KiB~?
\end{enumerate}

\medskip\noindent\textbf{Partie III --- Quand elles ne tiennent pas.}
\begin{enumerate}[resume]
\item Le desk veut 300 instruments. Quelle fraction des mises à jour manquera le deuxième niveau si les instruments sont mis à jour
uniformément au hasard, et si le deuxième niveau contient les plus récents~?
\item Si une mise à jour effectue trois lectures dépendantes, quel est son temps mémoire moyen avec ce taux de défaut, en prenant les
latences du deuxième niveau et de la mémoire de la partie~I~?
\item Qu'est-ce que les \omterm{def:ll:memory-hierarchy-and-caches:tlb}{grandes pages} changeraient, et pourquoi particulièrement ici~?
\item Citer deux changements de disposition qui réduisent les données chaudes.
\end{enumerate}

\medskip\noindent\textbf{Partie IV --- Le verdict.}
\begin{enumerate}[resume]
\item Énoncer le \emph{résultat nommé}~: la latence de chaque niveau sur ce portable et le nombre d'instruments dont le carnet
tient dans le deuxième niveau.
\item Pourquoi 32~octets est-il une bonne taille pour un niveau, et 48 une mauvaise~?
\item L'échelle doit-elle stocker les niveaux sous forme de structures à quatre champs ou de quatre tableaux d'un seul champ~? Quand chacune l'emporte-t-elle~?
\item Que mesure le parcours que masque un benchmark séquentiel~?
\item Pourquoi le benchmark doit-il visiter chaque ligne une fois par tour~?
\item Quel est le coût, par mise à jour, d'un compteur en \omterm{def:ll:memory-hierarchy-and-caches:false}{faux partage} dans le carnet~?
\item Comment confirmeriez-vous, sur l'hôte de production, que les carnets restent dans le deuxième niveau~?
\item En une phrase~: qu'est-ce qui détermine la vitesse d'un carnet d'ordres~?
\end{enumerate}
\end{problem}

\section{Questions d'entretien}

\begin{interviewq}[$\star$ \iqroles{développeur}]\label{iq:ll:memory-hierarchy-and-caches:1}
Pourquoi parcourir un \texttt{std::vector} est-il plus rapide que parcourir une \texttt{std::list} contenant les mêmes éléments~?
\end{interviewq}

\begin{interviewq}[$\star\star$ \iqroles{développeur}]\label{iq:ll:memory-hierarchy-and-caches:2}
Qu'est-ce que le \omterm{def:ll:memory-hierarchy-and-caches:false}{faux partage}~? Comment le détecteriez-vous et comment le corrigeriez-vous~?
\end{interviewq}

\begin{interviewq}[$\star\star$ \iqroles{développeur}]\label{iq:ll:memory-hierarchy-and-caches:3}
Qu'est-ce qu'un TLB, et pourquoi les systèmes de trading utilisent-ils des \omterm{def:ll:memory-hierarchy-and-caches:tlb}{grandes pages}~?
\end{interviewq}

\begin{interviewq}[$\star\star$ \iqroles{développeur}]\label{iq:ll:memory-hierarchy-and-caches:4}
Comment mesureriez-vous la latence de chaque niveau de la hiérarchie de caches d'une machine~?
\end{interviewq}

\begin{interviewq}[$\star\star$ \iqroles{développeur}]\label{iq:ll:memory-hierarchy-and-caches:5}
Deux tableaux de 4\,096 doubles sont parcourus élément par élément en même temps, et la boucle est étonnamment lente. Que peut-il
se passer~?
\end{interviewq}

\begin{interviewq}[$\star\star\star$ \iqroles{développeur}]\label{iq:ll:memory-hierarchy-and-caches:6}
Concevez la disposition mémoire d'une table de correspondance identifiant d'ordre vers ordre pour un flux comptant un million d'ordres vivants, avec la latence p99 la plus basse possible.
\end{interviewq}
